\chapter{Limitations and threats to validity}
\label{chap:10-limites}

This chapter gathers the limitations of the work, so that no reader has to hunt for them. It sorts
them into the four families of threats of \textcite{shadish2002experimental}. Internal validity asks
whether the measured effect really comes from the presumed cause. External validity asks to which
populations, tasks and conditions the result extends. Construct validity asks whether the measure
captures the notion it claims to capture. Statistical conclusion validity asks whether the tests and
intervals license the conclusion drawn. Two sections close the chapter: the experiments not yet
finished, then what a new complete training run on another machine would test.

\section{Internal validity}
\label{sec:val-interne}

\subsection{The survey readers are agentic AI models}

Every paper was read by AI agents, copies of one model (\texttt{claude-opus-5-5}). Two
readers finished on a neighbouring model (\texttt{claude-opus-5}) for thirteen papers, after a
response was stopped by a content filter of the provider (survey deviation 28). No human reread the
papers blind: the rereading planned by deviation 13 was dropped before it began (deviation 23). The
one of deviation 35, thirty papers of both years drawn at random, still awaits its reader
(\cref{sec:sm-doubles}). Five safeguards reduce the risk without removing it.

\begin{enumerate}
  \item \textbf{The mandatory quotation.} Every \enquote{yes} decision rests on an exact quotation
        and its page. Anyone can check every decision against the PDF. The script
        \file{study/check\_codes.py} checks mechanically that every quotation appears in the text.
  \item \textbf{The quotation audit.} After the first percentages, we looked for the quotations that
        the mechanical check did not find (deviation 15, file \file{study/v1\_quotes\_audit.md}). Of
        89 quotations, none is invented: 83 appear word for word in the paper, one is a faithful
        paraphrase, 5 give values that contradict the rules.
  \item \textbf{The blind double readings.} Thirty papers of the survey were reread by a second agent
        that did not see the first reading: the first ten of the first 2026 draw (readers G and G2),
        five 2018 papers and five 2026 papers (R2), ten papers of the yearly trajectory (R3). On the
        ten reread papers of the first draw, the two readers agree on 127 decisions out of 130
        (98\%). On the four verdict criteria, the thirty papers give thirty agreements, a Wilson
        interval of $[0.89, 1]$. Over the thirty papers pooled, every criterion exceeds 80\%
        agreement. The least reliable are the origin of the baselines (Y5) and the exclusions (Y6),
        25 agreements out of 30 each, with a Cohen's kappa of 0.65 and 0.62 whose bootstrap interval
        goes below 0.35 (\file{study/survey\_extra.json}, deviation 30 (a)). On the first draw, the
        performance measure (Y1) reached only 60\% agreement, before its rules were made more
        precise. The ten papers of the trajectory agree at 80\% or more on each of the 23 criteria
        (\cref{tab:accord}). Agreement on 2018 papers rests on the five papers of R2.
  \item \textbf{The search against the \enquote{no} decisions.} A \enquote{no} rests on no
        quotation. Every \enquote{no} on the verdict criteria was confronted with the text of the
        paper through keywords fixed in advance. Each excerpt found was then judged by an arbiter that
        saw neither the paper nor the year (deviation 27). Of 377 \enquote{no} decisions, 10 become
        \enquote{yes} (2 in 2018, 8 in 2026) and no verdict changes. If the 24 decisions judged
        doubtful are flipped as well, uncertainty (X2) and sensitivity (Y9) end up just above the
        threshold ($p = 0.07$ and $p = 0.054$).
  \item \textbf{Removing the neighbouring model.} Without the thirteen papers partly read by
        \texttt{claude-opus-5}, the verdicts hold, except reported uncertainty (X2), which no longer
        passes the Holm correction ($p = 0.07$, deviation 29).
\end{enumerate}

These safeguards leave one threat whole. The two readers of a double reading are copies of the same
model: their agreement measures the stability of that model, which differs from the agreement of two
independent readers. A systematic bias of the model, shared by all its copies, escapes the double
readings. The exact instructions, the model and the dates of every reader are published
(\file{study/readers/}), so that another reader can redo the reading. A reader from another family
of models, \texttt{gemma3:12b}, is rereading the 98 experimental papers of the main comparison
(survey deviation 33). Its reading has not finished yet (\cref{sec:val-en-cours}).

\subsection{Readers and year}

No main reader read both years. The 2018 papers were read by F, F2 and H1 to H4, the classical RL
papers of 2026 by other agents (\file{study/survey\_extra.json}, key \texttt{readers\_year}). All are
copies of the same model, but each session has its own instructions, date and version of the grid.
In the main comparison, an effect of the set of readers therefore cannot be told apart from an effect
of the year. Removing one reader at a time (deviation 21) does not detect it, since the effect would
come from the set as a whole.

Deviation 30 (d) measures what the data allow. Within one year, readers sometimes differ more than
chance would explain: in 2018, the share of Q5 goes from 0\% (H3, 0 papers out of 8) to 55\% (F, 6
out of 11), with $p = 0.049$ by permuting the readers. In 2026, that of X7 goes from 17\% to 83\%,
with $p = 0.027$. These gaps mix the reader with the batch of papers: F read the twelve papers of the
calibration, H1 to H4 the added papers. For the verdict criteria, only the second reader R2 read both
years. On its five papers per year, it finds, for all four criteria, the same differences between
years as the main readers. Five papers per year are still too few to close the question. Criterion
X8 closes it for itself alone, since each of its six main readers read both years (deviation 32).

Deviation 36 separates the year from the readers for the four verdict criteria. Eight fresh readers,
T1 to T8, reread the 98 papers, each of them on both years. A model with a random reader effect then measures
the year effect for a given reader. The verdicts on Q5, Y9 and X7 hold. The one on X2, already the most
fragile, keeps its direction without being established under the rule fixed in advance. The reader
standard deviation is zero or indistinguishable from zero (\cref{sec:sr-annees}). Two limits are left.
Readers T1 to T8 see the year in every PDF. They also come from the same model family as the original
readers: a bias of that family that varied with the year would affect both readings alike, so the
90\% to 100\% agreement between them does not rule it out. The reader from another family of models
(deviation 33) and the human rereading (deviation 35) target that bias. Neither has a result yet.

\subsection{Decisions taken after the results}

Several rules were fixed after results had been seen: the verdict rule of the survey
(deviation 20), the harmonisation of the 2018 calibration (deviation 15), the descriptive ranking
under modified physics (simulation deviation 11), the details of the GRPO bootstrap (deviation 13)
and the reset of the GRPO failure count (deviation 10). The twenty-one deviations born of the
technical review all come after the results. For eight simulation
deviations, that review had already computed a draft version of the numbers
(\cref{sec:ecarts-audit}). They are published as exploratory or corrective, next to the
preregistered results. \Cref{chap:09-ecarts} discusses each of these decisions. The order in which
the deviations before the repository became public were written rests on our word and on the commit dates
(\cref{sec:parole}).

\subsection{The simulation runs}

The reconstructed ledger counts 42 stopped attempts that the protocol did not declare
(\cref{sec:registres}). Every stop hit all the runs in progress at once, with no link to a score. For
SAC and CrossQ, the curves of the stopped attempts were overwritten: we cannot show that the relaunch
reproduces the same curve. The determinism of these two JAX programs has yet to be established
(deviation 32, not yet finished). The protocol announced Stable-Baselines3 2.6.0. Every training ran with
2.9.2, the same version for all runs. These deviations favour no program. They are stated here so
that the reader knows them.

\subsection{The GRPO strand}

Three defects of the GRPO strand affect the cause of the measured effects (deviation 18 of the GRPO
strand).
\begin{itemize}
  \item \textbf{The hardware changes with the random seed.} For Qwen, the two true-reward runs that
        show the gain ran on the GPU, the three that do not show it on the CPU. Five runs cannot
        separate these two causes (\cref{sec:llm-materiel}). The first three runs ran before the
        safeguards against a corrupted computation existed: a gradient that was corrupted but
        finite, clipped to 1, may have gone through them without a trace. The MPS computation is not
        deterministic from one attempt to the next. The Llama runs show it: at random seed 3, the
        true-reward run and the randomly-drawn-reward run start from the same model with the same
        randomness on the same GPU. Their greedy accuracy at step 0 is nevertheless 0.707 and 0.667.
  \item \textbf{Qwen ran with a repetition penalty of 1.1}, taken from its configuration file, which
        Llama does not have. The comparison between the two families therefore mixes the family with
        the decoding setting. Qwen's sampled answers also depart from the policy whose probabilities
        the training computes.
  \item \textbf{The pilot read the test set.} The task and the 160-token ceiling were chosen by
        looking at the test problems. No training saw that set.
\end{itemize}
The rerun of deviation 20 of the GRPO strand corrects all three defects (\cref{sec:val-en-cours}).

\section{External validity}
\label{sec:val-externe}

\subsection{The population of papers}

\begin{itemize}
  \item \textbf{The population has changed.} 27 of the 50 papers of the first 2026 draw deal with
        large language models, a family absent in 2018. Our verdicts apply to classical RL.
  \item \textbf{The title filter.} The regular expression applied to titles misses RL papers whose
        title does not say so, in 2018 as in 2026. We measured this for 2026 with a plan written in
        advance (deviation 24). Of 400 papers drawn outside the filter, 25 have an abstract that
        mentions RL, which comes to about 1{,}037 missed papers (from 708 to 1{,}503), 45\% of the
        whole. This is an upper bound, since the mechanical rule also counts false positives. Read
        with both grids, 20 of these papers resemble the first draw: code linked in 60\% of them (67\%
        in the draw), an uncertainty in 50\% (47\%). The 7 classical RL papers among them are too few
        to test our verdicts.
  \item \textbf{The title screen of the 2026 queue.} A classical RL paper whose title contains
        \enquote{reasoning}, \enquote{preference} or \enquote{math} could not enter the classical 2026
        sample. The 9 ICLR 2026 papers at the head of the queue, unreachable at the time of reading, were
        skipped.
  \item \textbf{The source of the PDFs.} The 2026 papers were read mostly in their arXiv version (45
        of the 50 classical RL papers), the 2018 papers in the conference proceedings. A link to code
        is often added to the arXiv version after acceptance: the rise of Q5 might have come from
        that. Survey deviation 31 measured it by rereading Q5 on the same kind of version in both
        years. With both years read in the proceedings, the rise is $+41$ points, $[+21, +56]$. With
        both years read in the latest arXiv version, it is $+34$ points, $[+15, +51]$. In the least
        favourable reading, it is still $+30$ points, $[+11, +47]$, with $p_{\text{Holm}} = 0.035$
        (\file{study/q5\_versions.json}, \cref{sec:sr-robustesse}). The asymmetry of versions is
        real, but it changes only four to six of the 48 papers of 2018: it explains part of the gap
        between our 2018 Q5 and Pineau's, not the verdict.
  \item \textbf{The trajectory without ICLR.} The years 2019 to 2024 cover only ICML and NeurIPS:
        OpenReview refuses automated downloads.
  \item \textbf{Pineau's draw.} Her 2018 sampling method is not described. Our comparisons with her
        figures assume a draw comparable to ours and serve as context: no verdict depends on them.
\end{itemize}

\subsection{The simulations}

The experiments are limited to the MuJoCo simulator and to four locomotion tasks
(\task{HalfCheetah}, \task{Hopper}, \task{Walker2d}, \task{Swimmer}), those of
\textcite{henderson2018matters}. They say nothing about image-based games, real robotics or offline
RL. In their preregistered form, E1 and B1 are limited to one task (\task{HalfCheetah}) and two
programs (PPO and TRPO). E1 redone over the sixteen (program, task) pairs extends the finding to
the four programs and the four tasks (deviation 22, \cref{sec:e1}). The four programs run with their
default hyperparameters, except in E6: another tuning might change the rankings. The modified
physics changes only two quantities, the masses and the friction.

\subsection{The budgets}

PPO trains for 1{,}024{,}000 steps, SAC and CrossQ for 200{,}000 (simulation deviations 1 and 14 (d)).
The equal-budget reading (deviation 12 (iii)) compares PPO at 200{,}000 steps, read off its
one-million-step curve: a PPO tuned for a budget of 200{,}000 steps might do better. The budget of
SAC and CrossQ, reduced to stay within the available compute, places them in the \enquote{little data}
regime. A longer budget might change their order. CrossQ's batch-renormalisation warm-up also lasts
100{,}000 updates: the critic spends half its budget in it, the actor its whole training, as
measured on the policy trained with random seed 0 on \task{Hopper} (deviation 28, \cref{tab:hyperparametres},
\cref{sec:programmes}).

\subsection{One platform}

Every run comes from one Mac with an Apple M2 Max chip, shared with other projects, under macOS and
for arm64. Another machine would give other curves for the same random seed
(\cref{sec:reentrainer}).

\subsection{The scale of the GRPO strand}

The GRPO strand is an illustration: two models of 0.5 and 1 billion parameters, one multiplication
task, 100 steps, groups of 8 answers, answers limited to 160 tokens, no KL term.
\textcite{shao2025spurious} work on 7-billion-parameter models specialised in mathematics, on
MATH-500, with far more compute. Our strand neither replicates nor refutes their experiment.

\section{Construct validity}
\label{sec:val-construit}

\begin{itemize}
  \item \textbf{Our grid reads 2018 differently from Pineau.} Our reading finds 55\% of 2018 papers
        giving all the hyperparameters of their method, where Pineau counted 90\%. The criteria
        share their wording but have different implicit decision rules. This is the reason for the
        same-grid verdict rule (deviation 20).
  \item \textbf{A satisfied criterion does not mean a practice done well.} Q5 is yes for a link to
        code, with no check that the code runs. Q8 is yes for an applied test, with no judgement of
        whether the test is appropriate. X2 is yes for any error bar, whereas the bar computed over
        the episodes of one training run is exactly the trap that experiment B1 shows
        (\cref{sec:b1}). Criterion X8, added after the results (deviation 32), asks whether the main
        uncertainty is computed over random seeds. It is in 35\% of the experimental papers of 2018
        and 42\% of those of 2026, a gap of $+7$ points, $[-12, +25]$, which the survey does not
        detect. The rise of X2 therefore measures the spread of the error bar more than that of the
        bar computed over random seeds (\cref{sec:sr-principale}). X8 was read after the results by
        eight other fresh agents, S1 to S8 (Claude Code, model \texttt{opus}), that had read no paper
        of the survey. They are distinct from readers T1 to T8 of deviation 36.
  \item \textbf{The performance measure (Y1)} long divided the readers: 60\% agreement on the first
        draw, before its rules were made more precise, 26 agreements out of 30 over all the double
        readings. The verdict on flaw F8 is read with this caveat.
  \item \textbf{The borderline cases of X3, X4 and X7} received no new harmonisation rule, so that
        both years are read with the same grid. A compute figure limited to inference time, or failed
        random seeds that were discarded, stay coded as the reader read them.
  \item \textbf{What a checklist declares.} The authors' answers to the NeurIPS checklist do not
        guarantee the practice they name: none of the 18 papers that answer yes to the question on
        statistical significance applies a test. The survey measures what the paper does, the
        checklist what it declares.
  \item \textbf{The score of a simulation run} is the mean return of ten deterministic episodes at
        the last step. Experiment B3 shows that another measure (best score, mean of the curve)
        changes the order of some pairs. Reporting the best point of the curve instead of the last
        inflates the interquartile mean of PPO by up to 48\% (PPO JAX on \task{Walker2d}, ten random
        seeds, deviation 29, \file{sim/analysis/curves.json}).
  \item \textbf{E3 compares two programs, not two implementations of one algorithm.} The two PPOs
        share their algorithm hyperparameters, not their default initialisation: Stable-Baselines3
        initialises every layer orthogonally, while sbx leaves its orthogonal option off and keeps
        Flax's default initialisation (\cref{sec:programmes}). Deviation 34 tests whether this
        difference explains E3.
  \item \textbf{The pass@1 of the GRPO strand} mixes correctness and concision, since a truncated
        answer counts as wrong (\cref{sec:llm-coupure}). For Qwen, the pass@1 gain under the true
        reward is explained entirely by the drop in truncated answers. The accuracy of finished
        answers does not rise (deviation 19 of the GRPO strand).
  \item \textbf{The cost and risk ratings} that tested the cost thesis come from three agentic AI
        models of three different sizes. Their agreement with one another does not guarantee that a
        researcher would rate the same way.
\end{itemize}

\section{Statistical conclusion validity}
\label{sec:val-stat}

\paragraph{The sample sizes of the survey} With 48 experimental papers from 2018 and 50 from 2026,
the survey sees only large changes. Deviation 30 (b) computes, for each criterion, the smallest rise
that Fisher's test detects four times out of five: about thirty points for most criteria, more at the
corrected level (\cref{tab:difference-detectable}, \cref{sec:sr-principale}). \enquote{No detectable
change} therefore does not mean \enquote{no change}. On Q3, Q6, Q7 and Y8, the survey sees rises of
11 to 22 points without being able to establish them. That of Q6 even has a Newcombe interval that
avoids zero before correction. The first draw had 14 classical RL papers, too few to detect a gap of
thirty points, hence the enlargement of deviation 12. The trajectory rests on 20 papers per year: a
share there has an uncertainty of about twenty points on either side.

\paragraph{Multiple comparisons} The survey compares 23 criteria in the main comparison, with a Holm
correction applied separately to the 11 criteria of grid v1 and the 12 of grid v2. The correction
protects the error rate of each family at the cost of power: a gap whose Newcombe interval avoids
zero without passing Holm, such as Q6 or Y12, is reported but makes no verdict. For E2, the 24
probability-of-improvement intervals are also given at the simultaneous 95\% level through the
Bonferroni correction (simulation deviation 18, \cref{tab:e2-bonferroni}). E3 tests four tasks: with
the Holm correction (deviation 23), only the difference on \task{Swimmer} at twenty random seeds
is still settled.

\paragraph{The fragile verdicts} The rise of reported uncertainty (X2) no longer passes the Holm
correction in two checks ($p = 0.07$ each time). Its observed power after correction is only 0.52
(deviation 30 (b)). This rise is the most fragile of the four changes. Code and compute pass the
threshold in all our checks, sensitivity to a hyperparameter in nearly all of them
(\cref{tab:robustesse}).

\paragraph{The \enquote{settled} rule of the simulations} A pair is established, \enquote{settled} in
the generated tables, when the 95\% interval of the probability of improvement excludes 0.5, or when the intervals of the two
interquartile means are disjoint (\cref{sec:regle}). Deviation 26 measures the actual false-positive
rate of this rule on two groups drawn from the same program, where no difference exists
(\cref{tab:regle-tranche}). The rule settles wrongly slightly more than once in
twenty at ten runs against ten, nearly once in ten at five against five. The condition on the
interquartile means adds almost nothing. A pair settled narrowly at ten random seeds, such as PPO
JAX against CrossQ on \task{Hopper} (\cref{sec:e7}), is read with this rate in mind.

\paragraph{The number of random seeds in the simulations} The main results rest on ten random seeds
per pair. Experiment E5 shows what this number allows: between two programs whose means differ by
less than 20\%, five random seeds recover the right order only 71 times in a hundred at the median,
62 times as measured on twenty random seeds. E5 judges the right order on the very sample it draws
from. Judged out of sample, on random seeds 10 to 19, the order drawn from random seeds 0 to 9 is
kept in 20 comparisons out of 24, but in only 5 out of 7 for the close comparisons. On those close
comparisons, five random seeds recover the out-of-sample order 65 times in a hundred on average
(deviation 25, \file{sim/analysis/e5\_holdout.json}). The power of E3 gives the same order of
magnitude: at ten random seeds per program, the test corrected for four tasks detects four times out
of five only a probability of improvement of at least 0.94, 0.81 at twenty (deviation 23). On the
tasks where the first program does not clearly beat the second, the rankings at ten random seeds
offer no better guarantee than these rates. The twenty random seeds of deviation 12 confirm or qualify each result.
\Cref{chap:06-experiences-resultats} states every divergence.

\paragraph{The number of random seeds in the GRPO strand} With five random seeds, the bootstrap has
only 126 distinct resamples, ten with three random seeds. The percentile intervals of the protocol
are too narrow: on five random seeds, they are only about 60\% as wide as a Student interval. The
exact sign test never goes below $p = 2/32 = 0.0625$ with five random seeds, so no gain within a
configuration can be established, even when all five random seeds go the same way (deviation 19 of
the GRPO strand). An interval that contains zero says that we cannot conclude. It does not say that
the effect is zero.

\paragraph{The bootstrap intervals} The rliable intervals depend on NumPy's global generator and on
the order of the calls. Their bounds move with the random seed of the bootstrap, at the level of Monte
Carlo noise (simulation deviation 13): a verdict whose bound grazes the threshold is fragile.

\section{The experiments not yet finished}
\label{sec:val-en-cours}

Six experiments are written in advance and not yet finished. The report gives
no result for any of them. Each answers a limitation of this chapter. The watchdog stopped
the simulation queue once on a memory alert, then the queue resumed.
\begin{itemize}
  \item \textbf{The determinism of SAC and CrossQ} (simulation deviation 32): three training runs
        repeated identically will tell whether these JAX programs give the same run again.
  \item \textbf{The Zoo recipe with $\gamma = 0.99$} (simulation deviation 33 (c)): ten training runs
        will tell whether the discount factor explains the agents that stay upright without moving
        forward.
  \item \textbf{The cause of the gap between the two PPOs} (simulation deviation 34): 160 training
        runs add power to E3 and swap the orthogonal initialisation.
  \item \textbf{A reader from another family of models} (survey deviation 33): \texttt{gemma3:12b}
        rereads the four verdict criteria on the 98 experimental papers.
  \item \textbf{The human rereading} (survey deviation 35): we will reread ourselves thirty papers of
        both years blind. The thirty sheets still await their reader.
  \item \textbf{The GRPO rerun} (deviation 20 of the GRPO strand): 40 runs on a single device, with
        no repetition penalty and a pilot kept out of the test set, wait for the GPU to become free.
\end{itemize}

\section{What a new training run on another machine would test}
\label{sec:reentrainer}

No training was redone on another machine. We expect a new complete training run on another
processor not to reproduce our curves. MuJoCo, JAX and PyTorch would not perform their roundings in
the same order there. Two trajectories that differ in the last digit of one observation then drift
apart step after step. The GRPO strand shows that one machine is not always enough: on the Mac's
GPU, the same random seed gave two different starting points (\cref{sec:val-interne}).

Such a training run would therefore test the statistics, not each curve. It would tell whether the
spread between random seeds, the rankings, the inflation due to selection and the losses under
modified physics stay of the same order with other draws. A conclusion that survived only our own
random seeds would say nothing about the algorithms. Continuous integration already runs the
equivalent test for the analyses: on Linux, the counts, the verdicts and the rankings come back
exactly, the floating-point values to within $10^{-12}$ in relative terms (\cref{sec:ci}).
